\section{Definición del Problema}
\label{sec:problema}

Se observa la superficie libre del agua y, cuando está disponible, el flujo, y
se desea recuperar la geometría del fondo que no se ha medido. Para que esa
recuperación sea posible se necesita un modelo físico que vincule lo observable
en la superficie con el fondo desconocido: las ecuaciones de aguas someras
(SWE). Esta sección deriva brevemente las SWE, enuncia las hipótesis bajo las
que son válidas, explica por qué su estructura permite la inversión, y fija la
notación empleada en el resto del documento.

\subsection{De Navier--Stokes a las ecuaciones de aguas someras}

El punto de partida son las ecuaciones de Navier--Stokes para un fluido
incompresible, es decir, la conservación de masa y de cantidad de movimiento
\cite{vreugdenhil1994, vallis2017}. Las SWE se obtienen integrándolas en la
vertical bajo las siguientes hipótesis:

\begin{enumerate}
  \item \textbf{Continuo e incompresibilidad.} El agua se trata como un medio
        continuo de densidad constante $\rho$.
  \item \textbf{Aguas someras (onda larga).} La escala vertical es mucho menor
        que la horizontal, $(H/L)^2 \ll 1$. La aceleración vertical es
        despreciable y la presión resulta hidrostática, $\pd{p}{z} = -\rho g$,
        de donde $p = \rho g\,(\eta - z)$.
  \item \textbf{Velocidad uniforme en la vertical.} El perfil vertical de
        velocidad se considera aproximadamente uniforme; todas las variables de
        las SWE son promedios en la profundidad.
  \item \textbf{Fondo impermeable y fijo.} No hay flujo a través del lecho y la
        batimetría no cambia en el tiempo, $\pd{\zb}{t} = 0$; en la superficie
        libre se impone la condición cinemática.
  \item \textbf{Cierre de fricción.} El esfuerzo de fondo se parametriza con la
        fórmula de Manning; la viscosidad molecular es despreciable frente a la
        turbulenta a números de Reynolds altos.
\end{enumerate}

Integrando la continuidad y la cantidad de movimiento entre el fondo
$\zb(x,y)$ y la superficie $\eta(x,y,t)$ mediante la regla de Leibniz, y
aplicando las condiciones de borde anteriores, se obtienen las SWE
bidimensionales en forma primitiva:

\begin{align}
  \pd{h}{t} + \pd{(hu)}{x} + \pd{(hv)}{y} &= 0,
    \label{eq:swe-cont}\\
  \pd{u}{t} + u\pd{u}{x} + v\pd{u}{y} &= -\,g\,\pd{(h+\zb)}{x} - S_{fx},
    \label{eq:swe-momx}\\
  \pd{v}{t} + u\pd{v}{x} + v\pd{v}{y} &= -\,g\,\pd{(h+\zb)}{y} - S_{fy},
    \label{eq:swe-momy}
\end{align}

donde $h$ es la profundidad, $(u,v)$ las velocidades promediadas en la
vertical, $\zb$ la batimetría y $S_{fx} = g\,n^2\,u\,\lVert\mathbf{u}\rVert /
h^{4/3}$ la pendiente de fricción de Manning (coeficiente $n$; análogo para
$S_{fy}$). La misma física admite formas equivalentes. En la forma
\emph{conservativa} las incógnitas son las cantidades conservadas
$(h, hu, hv)$ y las ecuaciones se escriben como la divergencia de unos flujos
más un término fuente que recoge la pendiente del fondo y la fricción; conserva
masa y momento de forma exacta, pero obliga a recuperar la velocidad como
$u = hu/h$, inestable cuando $h \to 0$ (zonas que se secan). La forma
\emph{primitiva-conservativa} multiplica el residuo primitivo por una matriz
que reescala automáticamente cada ecuación, uniendo la estabilidad de las
variables primitivas con la ponderación de la conservativa. La elección de la
forma del residuo empleada por el método se discute en la
Sección~\ref{sec:method}.

\subsection{Por qué la batimetría es recuperable}

El término de presión hidrostática introduce el acoplamiento
$-\,g\,\pd{(h+\zb)}{x} = -\,g\,\pd{\eta}{x}$, con $\eta = h + \zb$ la
elevación de la superficie libre. Ese término es el \emph{canal de información}
entre lo observable ($\eta$) y lo desconocido ($\zb$): una sola instantánea es
equifinal, pues infinitas particiones $(h,\zb)$ producen el mismo $\eta$, pero
la evolución temporal del flujo depende de $\zb$ a través de este acoplamiento,
lo que rompe la degeneración y vuelve tratable el problema inverso.

\subsection{Validez y límite de aplicabilidad}

Las SWE son válidas mientras $(H/L)^2 \ll 1$, equivalente a $k h \ll 1$ para
una perturbación de número de onda $k$: ondas largas, no dispersivas, que se
propagan a velocidad $c = \sqrt{g h}$ \cite{leveque2002, vallis2017}. Cuando
$k h$ se acerca a la unidad aparecen efectos dispersivos que las SWE no
capturan; este límite se manifiesta empíricamente en el Exp.~6
(Sección~\ref{sec:exp6}), con datos de canal de oleaje en régimen
$k h \approx 1.3$.

\subsection{Diferencias con Navier--Stokes}

Frente a las ecuaciones de Navier--Stokes tridimensionales, las SWE reducen el
problema de 3D a 2D; la hipótesis hidrostática elimina la ecuación vertical de
cantidad de movimiento; no se resuelve el perfil vertical de velocidad sino su
promedio; la turbulencia tridimensional no se resuelve, sino que se parametriza
mediante la fricción de fondo. El costo computacional es mucho menor, a cambio
de un dominio de validez acotado a flujos de aguas someras.

\subsection{Notación}

La Tabla~\ref{tab:notacion} resume la notación empleada en el resto del
documento.

\begin{table}[H]
\centering
\caption{Notación empleada en el documento.}
\label{tab:notacion}
\begin{tabular}{ll}
\toprule
Símbolo & Significado \\
\midrule
$x,\ y$              & Coordenadas horizontales \\
$t$                  & Tiempo \\
$h$                  & Profundidad del agua \\
$u,\ v$              & Velocidades promediadas en la vertical \\
$\zb$                & Batimetría (elevación del fondo), $\pd{\zb}{t}=0$ \\
$\eta = h + \zb$     & Elevación de la superficie libre \\
$g$                  & Aceleración de gravedad \\
$n$                  & Coeficiente de Manning \\
$q = h u$            & Caudal por unidad de ancho \\
$S_f$                & Pendiente de fricción \\
$Fr = u/\sqrt{g h}$  & Número de Froude \\
$c = \sqrt{g h}$     & Velocidad de onda larga \\
\bottomrule
\end{tabular}
\end{table}
